\section{Integral charge data for the compact examples}
\label{app:matrices}

\subsection[The X9 divisor reduction]{The $X_9$ divisor reduction}

In the ordered basis \eqref{eq:x9-divisor-basis}, the three local coordinate
rows are
\begin{equation}
 B_{\mathrm{red}}
 =\begin{pmatrix}
  1&0&-1&-1&0&0\\
 -1&1& 0& 1&0&0\\
  0&1&-1& 0&0&0
 \end{pmatrix},
 \qquad
 Q_{\mathcal D}=B_{\mathrm{red}}^{\mathsf T}.
 \label{eq:app-x9-b}
\end{equation}
The following integral vector combinations define a unimodular basis:
\begin{align}
 V_1&=-D_6,&
 V_2&=D_5,\nonumber\\
 S_1&=D_2+D_7,&
 S_2&=D_2+D_5+D_6,\nonumber\\
 S_3&=D_8,&
 \Phi&=E.
 \label{eq:app-x9-basis-change}
\end{align}
The determinant of the change-of-basis matrix is $-1$, and in this frame
\begin{equation}
 Q'=
 \begin{pmatrix}
 1&0&1\\
 0&1&1\\
 0&0&0\\
 0&0&0\\
 0&0&0\\
 0&0&0
 \end{pmatrix}.
\end{equation}
This proves both the effective matrix and its primitivity in the printed
toric lattice.

The full $E_2$ tangent scheme also contains the nilpotent coordinate $\alpha$.
Its row is the pairing of the divisors with a class of $K_{S_7}^{\perp}$
complementary to the root; the choice of complement is immaterial on the
reduced branch, where $\alpha=0$. In the divisor order
\eqref{eq:x9-divisor-basis}, the row is
\begin{equation}
 (0,-2,1,0,-2,0).
 \label{eq:app-x9-alpha}
\end{equation}
Some spectator combinations can therefore couple to $\alpha$ over dual
numbers. This does not add a finite Higgs mass because
$\alpha^2=\alpha\beta=0$ and $\alpha$ vanishes on the reduction.

\subsection[The X-circ toric and resolution input]{The $X^{\circ}$ toric and resolution input}
\label{app:xcirc-input}

Let $N_{\circ}\cong\ZZ^4$ be a fresh one-parameter-subgroup lattice and set
$\Pi=\operatorname{conv}\{r_1,\ldots,r_{26}\}\subset
(N_{\circ})_{\mathbb R}$. Thus $\Pi$ is the fan polytope: the ambient
fourfold is defined by the fan over its faces, and $X^{\circ}$ is a
$\Pi$-regular anticanonical hypersurface. In particular, $\Pi$ is not the
polar $\Delta^{\circ}\subset M_{\mathbb R}$ of
\eqref{eq:polar-convention}. The following table is the complete vertex list;
its order fixes all subsequent labels.
\begin{equation}
\begin{array}{c|c@{\qquad}c|c}
1&(-1,-1,-1,-1)&14&(0,-1,-1,0)\\
2&(-1,-1,-1,0)&15&(0,-1,0,1)\\
3&(-1,-1,0,-1)&16&(0,-1,0,0)\\
4&(-1,-1,0,1)&17&(0,1,-1,-1)\\
5&(-1,-1,1,0)&18&(0,1,-1,0)\\
6&(-1,-1,1,1)&19&(0,1,0,-1)\\
7&(-1,0,-1,-1)&20&(0,1,0,0)\\
8&(-1,0,-1,0)&21&(1,1,-1,-1)\\
9&(-1,0,0,-1)&22&(1,0,-1,0)\\
10&(-1,0,0,1)&23&(0,0,0,-1)\\
11&(-1,0,1,0)&24&(0,0,0,1)\\
12&(-1,0,1,1)&25&(0,0,1,0)\\
13&(0,-1,-1,-1)&26&(0,0,1,1)
\end{array}
 \label{eq:xcirc-vertices}
\end{equation}
The twenty-six unit-square faces, written by vertex index, are
\begin{equation}
\begin{array}{c|c@{\quad}c|c}
N_1&\{1,2,7,8\}&N_{14}&\{5,6,25,26\}\\
N_2&\{1,2,13,14\}&N_{15}&\{7,8,17,18\}\\
N_3&\{1,3,7,9\}&N_{16}&\{7,9,17,19\}\\
N_4&\{2,4,8,10\}&N_{17}&\{10,12,18,20\}\\
N_5&\{2,4,14,15\}&N_{18}&\{10,12,24,26\}\\
N_6&\{3,5,9,11\}&N_{19}&\{11,12,19,20\}\\
N_7&\{3,5,13,16\}&N_{20}&\{11,12,25,26\}\\
N_8&\{3,5,23,25\}&N_{21}&\{13,14,15,16\}\\
N_9&\{3,9,19,23\}&N_{22}&\{13,16,23,25\}\\
N_{10}&\{4,6,10,12\}&N_{23}&\{15,16,25,26\}\\
N_{11}&\{4,10,15,24\}&N_{24}&\{17,18,19,20\}\\
N_{12}&\{5,6,11,12\}&N_{25}&\{18,20,24,26\}\\
N_{13}&\{5,6,15,16\}&N_{26}&\{19,20,25,26\}
\end{array}
 \label{eq:xcirc-node-faces}
\end{equation}
Each set in \eqref{eq:xcirc-node-faces} denotes the corresponding vertices
$r_i$. The remaining singular faces are
\begin{equation}
\begin{array}{c|c|c}
D_{6,1}&\{r_1,r_2,r_3,r_4,r_5,r_6\}&r_{27}=(-1,-1,0,0)\\
D_{6,2}&\{r_7,r_8,r_9,r_{10},r_{11},r_{12}\}&r_{28}=(-1,0,0,0)\\
D_{7,1}&\{r_1,r_7,r_{13},r_{17},r_{21}\}&r_{29}=(0,0,-1,-1)\\
D_{7,2}&\{r_2,r_8,r_{14},r_{18},r_{22}\}&r_{30}=(0,0,-1,0),
\end{array}
 \label{eq:xcirc-divisorial-faces}
\end{equation}
where the last column is the unique interior lattice point. These face data
give precisely the singularity inventory in \eqref{eq:xcirc-content}.

For definiteness, one smooth projective crepant subdivision is fixed as
follows. On the thirty boundary rays assign heights
\begin{equation}
 b_i=2^{26-i}\quad(1\leq i\leq26),
 \qquad
 b_i=-2^{30}+2^{30-i}\quad(27\leq i\leq30).
 \label{eq:xcirc-resolution-heights}
\end{equation}
Take the lower-facet triangulation on every facet and cone it to the origin.
The induced triangulations agree on common faces, every maximal cone is
unimodular, and the four divisorial faces are star-triangulated at
$r_{27},\ldots,r_{30}$. This fixes the exceptional curves and surfaces used
below.

Principal-divisor relations eliminate the rays $15,22,24,26$. Thus the
rank-$26$ divisor basis is indexed by
\begin{equation}
 \mathcal I=
 \{1,2,\ldots,14,16,17,18,19,20,21,23,25,27,28,29,30\}.
 \label{eq:app-xcirc-index}
\end{equation}
For a local coordinate $a$, write its charge row as
$b_a=\sum_{I\in\mathcal I}B_{aI}D_I$. The complete sparse presentation of
$B$ is
\begingroup
\scriptsize
\begin{align*}
b_{N_1}&=D_1-D_2-D_7+D_8,&
b_{N_2}&=D_1-D_2-D_{13}+D_{14},\\
b_{N_3}&=D_1-D_3-D_7+D_9,&
b_{N_4}&=D_2-D_4-D_8+D_{10},\\
b_{N_5}&=D_2-D_4-D_{14},&
b_{N_6}&=D_3-D_5-D_9+D_{11},\\
b_{N_7}&=D_3-D_5-D_{13}+D_{16},&
b_{N_8}&=D_3-D_5-D_{23}+D_{25},\\
b_{N_9}&=D_3-D_9+D_{19}-D_{23},&
b_{N_{10}}&=D_4-D_6-D_{10}+D_{12},\\
b_{N_{11}}&=D_4-D_{10},&
b_{N_{12}}&=D_5-D_6-D_{11}+D_{12},\\
b_{N_{13}}&=D_5-D_6-D_{16},&
b_{N_{14}}&=D_5-D_6-D_{25},\\
b_{N_{15}}&=D_7-D_8-D_{17}+D_{18},&
b_{N_{16}}&=D_7-D_9-D_{17}+D_{19},\\
b_{N_{17}}&=D_{10}-D_{12}-D_{18}+D_{20},&
b_{N_{18}}&=D_{10}-D_{12},\\
b_{N_{19}}&=D_{11}-D_{12}-D_{19}+D_{20},&
b_{N_{20}}&=D_{11}-D_{12}-D_{25},\\
b_{N_{21}}&=D_{13}-D_{14}-D_{16},&
b_{N_{22}}&=D_{13}-D_{16}-D_{23}+D_{25},\\
b_{N_{23}}&=-D_{16}+D_{25},&
b_{N_{24}}&=D_{17}-D_{18}-D_{19}+D_{20},\\
b_{N_{25}}&=D_{18}-D_{20},&
b_{N_{26}}&=D_{19}-D_{20}-D_{25},\\
b_{D_{6,1},s_1}&=D_1-D_2-D_5+D_6,&
b_{D_{6,1},s_2}&=D_2-D_3-D_4+D_5,\\
b_{D_{6,1},s_3}&=D_1-D_2-D_3+D_4+D_5-D_6,&
b_{D_{6,2},s_1}&=D_7-D_8-D_{11}+D_{12},\\
b_{D_{6,2},s_2}&=D_8-D_9-D_{10}+D_{11},&
b_{D_{6,2},s_3}&=D_7-D_8-D_9+D_{10}+D_{11}-D_{12},\\
b_{D_{7,1},\alpha}&=-2D_1+2D_7+D_{13}-D_{17},&
b_{D_{7,1},\beta}&=-D_1+D_{13}+D_{17}-D_{21},\\
b_{D_{7,2},\alpha}&=2D_2-D_8-2D_{14},&
b_{D_{7,2},\beta}&=D_8-D_{14}-D_{18}.
\end{align*}
\endgroup
This is a $36\times26$ matrix of rank $21$, with twenty-one nonzero Smith
invariants, all equal to one. Deleting the coordinates absent on a chosen
$E_3$ component gives the four profile matrices in
Table~\ref{tab:xcirc-profiles}.

\subsection[Primitive vectors in X-circ]{Primitive vectors in $X^{\circ}$}

If $B$ denotes the $36\times26$ local-coordinate-by-divisor matrix, the five
primitive vectors producing \eqref{eq:xcirc-isolating-moments} are
\begin{align}
 v_9={}&\sum_{i=1}^{12}D_i+D_{17}+D_{18}+D_{19}+D_{20},\nonumber\\
 v_{11}={}&-\sum_{i=7}^{12}D_i-D_{17}-D_{18}-D_{19}-D_{20}-D_{21},\nonumber\\
 v_{\beta_1}={}&-D_{21},\nonumber\\
 v_{\beta_2}={}&\sum_{i=1}^{12}D_i-D_{21},\nonumber\\
 v_{A_1}={}&D_1+D_3+D_5+D_7+D_9+D_{11}+D_{13}+D_{16}
 +D_{17}+D_{19}+D_{21}+D_{23}+D_{25}.
 \label{eq:app-xcirc-vectors}
\end{align}
Direct multiplication gives
\begin{align}
 Bv_9&=e_{u_9}-e_{\alpha_1}+e_{\alpha_2},&
 Bv_{11}&=e_{u_{11}}-e_{\alpha_1}+e_{\alpha_2},\nonumber\\
 Bv_{\beta_1}&=e_{\beta_1},&
 Bv_{\beta_2}&=e_{\alpha_2}+e_{\beta_2},\nonumber\\
 Bv_{A_1}&=e_{L_1}+e_{L_2}.
 \label{eq:app-xcirc-products}
\end{align}
The complete sparse matrix is included in the reproducibility source described
in Appendix~\ref{app:reproducibility}. Its Smith form has $21$ nonzero unit
invariants and five zero directions. After branch restriction, the four
profiles have the ranks printed in Table~\ref{tab:xcirc-profiles}.
